\section{Introduction}


\begin{figure}[!t]
    \centering
    \includegraphics[width=\linewidth]{figures/Diff2Lip-Teaserv5.pdf}
    \caption{\textbf{Top:} Our \ours~~approach uses an audio-conditioned diffusion model to generate lip-synchronized videos. (Here $q$ denotes the forward diffusion process and $p_\theta$ is the learned reverse diffusion process.) \textbf{Bottom:} On zooming in to the mouth region it can be seen that our method generates high-quality video frames without suffering from identity loss.}
    \label{fig:teaser}
    \vspace{-0.2in}
\end{figure} 



Oscar-winning director Bong Joon Ho famously pointed out that subtitles act as a barrier between a foreign audience not adept in the language and their ability to fully enjoy amazing movies~\cite{subtitle-quote}, as the viewer needs to focus on both watching and reading. A rarely explored alternative, multiple-language version movie (MLV), where the same film is shot in multiple languages in parallel, is naturally much more expensive~\cite{mlv}. While dubbing is a popular compromise solution, it can feel unnatural due to the lack of synchronization between speech and actors' video. 
As a cheaper alternative, lip-synchronization (lip-sync) aims to generate the mouth region of the human face such that the lips correspond to a different speech audio. Its applications beyond movies include education, virtual avatars, and video conferencing.  
Ideally, lip-sync should support any identity and audio from unseen sources (in-the-wild setting). 
This setting brings up the challenges of preserving the actors’ identity, pose, emotions, and visual quality while maintaining a realistic lip-sync.

One of the earliest lip-sync methods, Video Rewrite~\cite{video-rewrite}, had a purpose-built solution by mapping phonemes to mouth shapes and then blending them onto the target video. Modern techniques have more general solutions but suffer certain limitations. 
For instance, PC-AVS~\cite{pcavs} and GC-AVT~\cite{gcavt}, disentangle pose and expression respectively but fail to preserve identity (Fig.~\ref{fig:teaser} bottom), have worse visual quality, and have border inconsistencies (while putting the generated heads back to the scene). On the other hand, works that target a specific identity, such as SynthesizingObama~\cite{synthesizing-obama}, require video/identity-specific training. Other methods which rely on extracting intermediate representations, e.g., landmarks in MakeItTalk~\cite{makeittalk}, have to deal with estimation errors in these representations. Finally, approaches that can generalize on in-the-wild lip-sync settings pose it as an inpainting task, where the mouth region is masked and then generated according to the audio. Examples include Wav2Lip~\cite{wav2lip}, which achieves good lip-sync but at the cost of poor visual quality (see Fig.~\ref{fig:teaser} bottom), and AV-CAT~\cite{avcat}, which has a multistage pipeline but does not capture finer details. In this paper, we introduce \ours, an inpainting style approach that solves the lip-sync task using diffusion models, which addresses most of these shortcomings and achieves visually superior lip-sync results. 


























 





We propose an audio-conditioned diffusion model to solve the task of lip-sync (Fig.~\ref{fig:teaser} top). 
Diffusion models~\cite{ddpm} are likelihood-based models that can generate astonishing results in high variation datasets~(e.g., ImageNet~\cite{imagenet}), that GANs~\cite{gan} cannot match. To generalize in-the-wild, we pose the problem as a conditional diffusion model based inpainting task~\cite{palette}. \ours  ~takes three inputs: a masked input frame, a reference frame, and an audio frame, and outputs the lip-synced mouth region. \ours  ~leverages (1) the masked input frame to get the pose context; (2) the reference frame to get the identity and mouth region textures; (3) the audio frame to drive the lip shape.
Using an audio+image conditioned diffusion model, \ours  ~maintains a fine balance between all these contextual input information, avoiding lip-sync problems (e.g. identity loss, reference copying, inaccurate lip shape). \ours  ~optimizes three losses: a reconstruction loss to guide synthesis; a sync-expert loss~\cite{wav2lip} to enforce synchronization; and a sequential adversarial loss to enforce inter-frame continuity.
\ours  ~generates high image quality without identity loss or generalizability issues as shown in Fig.~\ref{fig:teaser} bottom.



 






We evaluate our work on commonly used benchmarks of Voxceleb2~\cite{voxceleb2} and LRW~\cite{lrw} datasets for the tasks of reconstruction and cross generation (see section~\ref{sec:exp}). We compare against popular methods used for lip-sync like Wav2Lip~\cite{wav2lip} and PC-AVS~\cite{pcavs}. 
Extensive evaluations show that \ours ~outperform existing methods in terms of image fidelity while having comparable synchronization.


The following are the contributions of this work:
\begin{itemize}[noitemsep,topsep=0pt]
    \item  
    We propose a novel diffusion model based approach for audio-conditioned image generation.
    \item Using frame-wise and sequential losses we are able to successfully generate high quality lip-sync.
    \item We show that the use of a sequential adversarial loss makes frame-wise video generation more stable for diffusion models across frames.
    \item Extensive evaluations validate that our generations outperform existing methods in FID metric and MOS of the users showing the effectiveness of \ours. 
\end{itemize}

